\documentclass{article}

\usepackage{amsmath}
\usepackage{amssymb}
\usepackage{xcolor}
\usepackage{tikz}

\usepackage[utf8]{inputenc}
\usepackage[T1]{fontenc}

\usepackage[polish]{babel}

\usepackage{listings}
\lstset{
    language=SQL,
    inputencoding=utf8x,
    extendedchars=\true,
    literate={ą}{{\k{a}}}1
             {Ą}{{\k{A}}}1
             {ę}{{\k{e}}}1
             {Ę}{{\k{E}}}1
             {ó}{{\'o}}1
             {Ó}{{\'O}}1
             {ś}{{\'s}}1
             {Ś}{{\'S}}1
             {ł}{{\l{}}}1
             {Ł}{{\L{}}}1
             {ż}{{\.z}}1
             {Ż}{{\.Z}}1
             {ź}{{\'z}}1
             {Ź}{{\'Z}}1
             {ć}{{\'c}}1
             {Ć}{{\'C}}1
             {ń}{{\'n}}1
             {Ń}{{\'N}}1
}

\title{Czym jest $e^{i\phi}$?}
\date{\today}

\begin{document}
\maketitle

%%%%%%%%%%%%%%%%%%%%%%%%%%%%%%%%%%%%%%%%%%%%%%%%%%%%%%%%%%%%%%%
\section{Intuicja geometryczna}
Mnożenie przez liczbę zespoloną $1 + i\frac{\phi}{n}$, gdzie $n \rightarrow \infty$,
jest interpretowane jako obrót wektora o nieskończenie mały kąt $\frac{\phi}{n}$.

Jeśli wykonamy $n \rightarrow \infty$ takich obrotów, to ostatecznie
obrócimy wektor o kąt $\phi$, czyli przemnożymy przez liczbę
zespoloną $\cos{\phi} + i \sin{\phi}$.

Wiemy zatem, że granica ciągu $\left( 1 + i\frac{\phi}{n} \right)^{n}$
jest równa $\cos{\phi} + i \sin{\phi}$. A więc liczymy granicę:
\[
    \lim_{n \rightarrow \infty} \left( 1 + i\frac{\phi}{n} \right)^{n} = a_{n}
    \implies
    \log a_{n} = n\log \left( 1 + i\frac{\phi}{n} \right)
\]
Korzystamy z własności:
\[
    \log(1 + i\frac{\phi}{n}) = \frac{i\phi}{n} - \frac{(i\phi)^2}{2n^2} + \frac{(i\phi)^3}{3n^3} + \cdots
\]
Wówczas:
\[
    \lim_{n \rightarrow \infty}\log a_{n} = \lim_{n \rightarrow \infty} n\left( \frac{i\phi}{n} - \frac{(i\phi)^2}{2n^2} + \frac{(i\phi)^3}{3n^3} + \cdots \right) = i\phi
\]
Pozbywając się logarytmu otrzymujemy:
\[
    \lim_{n \rightarrow \infty} a_{n} = \lim_{n \rightarrow \infty} \left( 1 + i\frac{\phi}{n} \right)^{n}  = e^{i\phi}
\]

Zatem mamy równość:
\[
    e^{i\phi} = \cos{\phi} + i \sin{\phi}
\]


%%%%%%%%%%%%%%%%%%%%%%%%%%%%%%%%%%%%%%%%%%%%%%%%%%%%%%%%%%%%%%%
\section{Rozwinięcie w szereg Taylora}
Rozwinięcie eksponenty w szereg Taylora ma postać:
\[
    e^{t} = \sum_{k=0}^{\infty}{\frac{t^k}{k!}}.
\]
W przypadku $e^{i\phi}$ mamy:
\[
    \begin{array}{l}
    e^{i\phi} = \sum_{k=0}^{\infty}{\frac{(i\phi)^n}{k!}} 
    =
    1 + i\phi - \frac{\phi^2}{2!} - \frac{i\phi^3}{3!} + \frac{\phi^4}{4!} + \cdots =
    \\
    \underbrace{1 - \frac{\phi^2}{2!} + \frac{\phi^4}{4!} + \cdots}_{\cos{\phi}}
    + i \left( \underbrace{\phi - \frac{\phi^3}{3!} + \frac{\phi^5}{5!}}_{\sin{\phi} + \cdots} \right) =
    \cos{\phi} + i\sin{\phi}
    \end{array}
\]
Można zauważyć, że części rzeczywiste i urojone grupują się w
szeregi Taylora odpowiednio dla $\cos{\phi}$ i $\sin{\phi}$.

%%%%%%%%%%%%%%%%%%%%%%%%%%%%%%%%%%%%%%%%%%%%%%%%%%%%%%%%%%%%%%%
% \section{Równanie różniczkowe}


\end{document}